\documentclass[protokoll]{dvd}

\KOMAoptions{
	headwidth = 18cm,
	footwidth = 18cm,
}

\begin{document}

\title{Divisionsstämmans möte 2}
\subtitle{2023}
\author{Divisionsstämman}
\date{2023-05-24}

\textbf{Datum:} \csname @date\endcsname\\
\textbf{Tid:} 17:17\\
\textbf{Plats:} Monaden (ev. EA)\\
\textbf{Styrelsemedlemmar:}
\begin{närvarande_förtroendevalda}
\förtroendevald{Ordförande}{Samuel Hammersberg}{}
\förtroendevald{Vice ordförande}{\emph{Vakant}}{}
\förtroendevald{SAMO}{Tekla Siesjö}{}
\förtroendevald{Kassör}{Lukas Gartman}{}
\förtroendevald{Sekreterare}{\emph{Vakant}}{}
\end{närvarande_förtroendevalda}

% \textbf{Närvarande övriga medlemmar:}
% \begin{närvarande_medlemmar}
% \end{närvarande_medlemmar}

\section{Öppnande av möte}
Mötet beräknas öppnas av Samuel Hammersberg 17.17

\section{Formalia}

\subsection{Divisionsstämmans beslutbarhet}

6 kap. i stadgan definierar regler Divisionstämman.

Den första februari kallade styrelsen till divisionsstämma genom att skriva i discordservern \emph{MonadenDV};

Tyvärr kunde inte ett utskick ske över mail då divisionen har förlorat
access till studenternas maillistor då när Göteborgs Universitet
bytte över till Outlook, ändrades de även reglerna angående mailutskick
till maillistor. Styrelsen jobbar med att fixa access igen.

Dessa möteshandlingar ska ha skickats ut under måndag den 22de maj.

\subsubsection*{Förslag till beslut}

\begin{attsatser}
    \item divisionsstämman har uppnått kraven i stadgan för att få hålla möte, och är därmed beslutbar.
\end{attsatser}

\subsection{Fastställande av mötesschema}

För att divisionsstämman ska kunna fatta ett beslut eller protokollföra en diskussion behöver punkten i mötesschemat där stämman ska fatta beslut vara inlagd eller föras in i mötesschemat senast vid den här punkten.

\subsubsection*{Förslag till beslut}

\begin{attsatser}
    \item införa fastställande av röstlängden som punkt 2.1 i dagordningen.
    \item införa val av DVarms ordförande för verksamhetsåret 2023 under beslutsärenden.
    \item införa punkten "Planering för integreringen av Lindholmarna" under diskussionspunkter.
    \item införa punkten "Hur sköts Monaden" under diskussionspunkter.
    \item införa "avslutande av möte" i slutet av dagordningen
\end{attsatser}

\subsection{Val av mötesordförande}

Mötesordförande har till uppgift att leda Divisionsstämmans sammankomst.
Hen ansvarar för att mötesformalia sköts.

\subsubsection*{Förslag till beslut}

\begin{attsatser}
    \item Samuel Hammersberg väljs till mötesordförande.
\end{attsatser}

\subsection{Val av vice mötesordförande}

Vice mötesordförande hjälper mötesordförande med att hålla talarlistan, och att alla får komma till tals.

\subsubsection*{Förslag till beslut}

\begin{attsatser}
    \item Lukas Gartman väljs till vice mötesordförande
\end{attsatser}

\subsection{Val av mötessekreterare}

Mötessekreteraren har till uppgift att anteckna diskussioner, beslut, och eventuella reservationer under mötet.

\subsubsection*{Förslag till beslut}

\begin{attsatser}
    \item Tekla Siesjö väljs till mötessekreterare.
\end{attsatser}

\subsection{Val av protokolljusterare}

Protokolljusterare har till uppgift att kontrollera att protokollet i slutändan reflekterar de faktiska besluten och diskussionerna som fördes under sammanträdet; samt agera rösträknare vid slutna omröstningar.
Utöver protokolljusterarna så ska mötesordförande och mötessekreteraren signera protokollet.
Vid Divisionsstämmans sammanträden ska det vara två justerare.
Mötesordförande och mötessekreteraren kan inte vara justerare.

\subsubsection*{Förslag till beslut}

\begin{attsatser}
    \item {\it inga förslag från styrelsen innan mötet}.
\end{attsatser}

\newpage

% \section{Rapporter}
% \subsection{Arbetsgrupp för logga arbete}
% Under förra stämman så startades det en arbetsgrupp styrd av Tim Persson
% med syfte att ta fram en ny logotyp för divisionen. Under kommande stämman
% kommer en presentation framföras av gruppen för att visa hur arbetet har gått.

% \section{Rapporter}
% 
% \subsection{Styrelsen}
% 
% \subsubsection*{Verksamhetsrapport}
% 
% Styrelsen har haft fyra möten sedan Divisionens stämma den 17 maj.
% Protokollen från mötena finns (snart) tillgängliga för allmänheten på styrelsens drive.
% Länk till styrdokument och protokoll: \verb|https://drive.google.com/drive/folders/1o_lLM7g7ph-xdxgut3E_BU0ywichWYTX?usp=sharing|
% 
% Efter en lång kamp, så har vi äntligen fått en rullstolsramp till monaden, men kampen är inte över än, då vi fortfarande behöver ett större kök.
% 
% Samuel har varit på institutionsråd och programråd. På programrådet så togs saker som introkursen, obligatoriska kurser och möjligheten om ett tillägg av en kurs inom open source licenser. Vi åt väldigt mycket fika.
% 
% Mottagningen har även gått bra, och alla kommittéer har gjort ett bra jobb! Verkar även som att de fortsätter att rulla på bra!
% 
% \newpage
% 
% \subsubsection*{Ekonomi}
% 
% Den gångna tiden har varit en ny epok för Datavetenskapsdivisionen i avseende finanser och ekonomi. Mottagningsbudgeten 2022 var i relation till tidigare mottagningsbudgetar, enorm(!). Av vår tilldelade Göta-budget, för mottagningen, om 30tkr har 14tkr spenderats. Av vår tilldelade institutions-budget om 78tkr har 71tkr spenderats. Det totala beloppet för mottagningen i skrivande stund är alltså 85tkr. Totalt för året har vi haft kostnader om 95tkr och inkomster om 92tkr. Med det sagt, så har vi utstående inbetalningar från Göta som gör att summan förväntas landa på plus för kalenderåret.
% 
% Det nya bokföringssystemet har varit en lärandeprocess men då vi inte har något bokföringskrav så har vi tämligen lösa regler för vad som "krävs" och vi kan därmed utefter behov skapa konton och resultatgrupper för att göra resultrapporter "berättande" för medlemmar. En struktur för bokföringen är nu, och kommer förbli, en levande mall där tydlighet gentemot medlemmarna är en stor del av anledningen till bokföringssystemet. Jag tror och hoppas att nästa kassör med enkelhet kommer kunna fortsätta utveckla strukturen.
% 
% Då detta är min sista divisionsstämma som styrelsemedlem och kassör, så önskar jag att tacka för mig.
% 
% mvh, Morgan
% 
% (Notera att siffrorna är levande stämmer för 2022-11-20)
% 
% Tillägg: utöver mottagningen är ekonomin oförändrad, speciellt då det inte har gått så många skolveckor sedan förra stämman.
% 
% \newpage
% 
% \subsection{Verksamhetsrapporter av kommittéer}
% 
% \subsubsection*{DVRK}
% Muntlig verksamhetsrapport av DVRKs ordförande Leo Mirzajanzadeh.
% 
% \subsubsection*{ConCats}
% Muntlig verksamhetsrapport av Tekla Siesjö.
% 
% \subsubsection*{Femme++}
% Muntlig verksamhetsrapport av Femme++s ordförande Emmie Berger.
% 
% \subsubsection*{DVArm}
% Muntlig verksamhetsrapport av DVArms ordförande Albin Otterhäll.
% 
% \subsubsection*{Mega6}
% Muntlig verksamhetsrapport av Petter Blomkvist.
% 
% \newpage
% 
\section{Beslutsärenden}

Enligt Stadgan måste ändringar av Stadgan röstas igenom på två av Divisionsstämmans varandra följande möten.
Om en beslutpunkt innehåller ``första läsningen' innebär det att det är första gången beslutet tas upp för omröstning.
Om en beslutspunkt innehåller ``andra läsningen'' innebär det att beslutspunkten har röstats igenom förra stämmomötet, och beslutet behöver bekräftas för att gå igenom.

\subsection*{Proposition: Förlänging av tiden på väggundersökning}


\section{Diskussioner}\label{sec:discussioner}
Detta är en en ny punkt som lades in förra året,
och styrelsen skulle vilja att detta blir återkommande punkt.

\section{Avslutande av möte}

Mötet beräknas avslutas 18.18

\stämmosignaturer

\end{document}
